% Resumo: 150 a 250 palavras (enunciado). Rever no fim, com os números do TESTE.
\begin{abstract}
Um projeto de verificação de factos precisa de priorizar notícias suspeitas sem as ler uma a uma.
Estudamos que características linguísticas distinguem notícias falsas de verdadeiras em português
e se um detetor generaliza para temas não vistos, usando o corpus Fake.Br (7\,200 notícias
emparelhadas por assunto). Construímos um pipeline completo com divisão de teste agrupada por par,
para evitar fugas de informação, e três tarefas: classificação, modelação de tópicos e extração de
entidades. Mostramos que o corpus tem um atalho forte: um classificador que usa apenas o comprimento
do texto obtém F1-macro de 0,944 em validação cruzada, perto dos 0,974 do melhor modelo TF-IDF.
Descobrimos ainda que a versão normalizada do corpus cria um segundo atalho: as notícias
verdadeiras acabam quase sempre a meio de uma frase, e o último carácter basta para 0,890. Numa versão
sem os dois atalhos, os n-gramas de caracteres obtêm 0,939 e o BERTimbau afinado 0,987 no teste, e a remoção
de \emph{stopwords} e o \emph{stemming} pioram os resultados, o que indica que o sinal está no
estilo e não no vocabulário temático. Os tópicos LDA estão mais associados
ao tema (NMI 0,104) do que à veracidade (NMI 0,047), e, segundo 100 notícias anotadas manualmente, as verdadeiras
contêm cerca de cinco vezes mais datas. Os modelos de estilo perdem menos desempenho em temas não vistos
do que os de palavras. Discutimos o viés de fonte do corpus (93\% das notícias falsas vêm de um
único sítio), as implicações éticas e um conjunto de recomendações para a organização.
\keywords{Notícias falsas \and Português \and Estilometria \and Modelação de tópicos \and Extração de informação}
\end{abstract}
